% !PRINTED 519
\weberpage{519}
If now $m$ is divisible by a still higher power of 2 than the second, then, when the exponent of 2 is odd, one can by § 105, (7) reduce all the characters to such as are of the form
\[
\left(\frac{A}{m''}\right), \ \left(\frac{-1}{A}\right)\left(\frac{A}{m''}\right)
\]
Put
\[
m = 2r^2 m'',
\]
where again by $r^2$ is understood the greatest square number dividing $\tfrac12 m$, and apply formula (8) or (18) according as $m'' \equiv 1$ or $\equiv 3\ (\mathrm{mod}\ 4)$. One then obtains the complete splitting of the class equation upon adjunction of the square roots of all the primes dividing $m$, with the exception of $\sqrt2$.

But if the exponent of the highest power of 2 dividing $m$ is an even number, then (18) also does not yet suffice for the complete decomposition of the class equation. Here again one can indeed split the function $H$ by the formulae (8), (18) through the adjunction of all the $\sqrt p$ and $\sqrt q$, but $\sqrt2$ is not thereby brought in, and so two genera always still remain contained in each such partial equation. In order to separate these also, we derive yet a third formula, on the assumption that $m$ is divisible by 8.

We set
\[
m = m'm'',
\]
\[
x = 1, \quad y = \tfrac12 m' - 2m'', \quad n = \left(\tfrac14 m' + m''\right)^2.
\]
\[
\sqrt{-i(\alpha+\beta\omega)} = e^{-\frac{\pi i}{4}}\left(\tfrac12\sqrt{m'} + i\sqrt{m''}\right),
\]
where again $m''$ is odd and divisible by no square. Here it follows from § 136:
\weq{19}{M = (-1)^{\frac{m}{8}} e^{\frac{\pi i}{4} A m''} \left(\frac{A}{m''}\right) \left(i^{-\frac{m''+1}{2}} \tfrac12\sqrt{m'} + i^{-\frac{m''-1}{2}}\sqrt{m''}\right)^3.}

After the decomposition according to the characters $\left(\dfrac{A}{m''}\right)$ has been settled by the formulae (8) and (18), it is now only a question of the decomposition according to the characters
\[
\left(\frac{-1}{A}\right), \ \left(\frac{2}{A}\right),
\]
that is, according to the behaviour of $A$ with respect to the modulus 8. It is therefore sufficient to set $m'' = 1$ in formula (19), whereby one obtains the following four values of $M$:
%CONT>>
